\documentclass[12pt,a4paper]{article}
\usepackage[utf8]{inputenc}
\usepackage[italian]{babel}
\usepackage{graphicx}
\usepackage{hyperref}
\usepackage{listings}
\usepackage{xcolor}
\usepackage{geometry}
\usepackage{fancyhdr}
\usepackage{amsmath}
\usepackage{amssymb}
\usepackage{multicol}

\geometry{margin=2.5cm}

\lstset{
    basicstyle=\ttfamily\small,
    keywordstyle=\color{blue},
    commentstyle=\color{gray},
    stringstyle=\color{red},
    showstringspaces=false,
    breaklines=true,
    frame=single,
    numbers=left,
    numberstyle=\tiny\color{gray}
}

\pagestyle{fancy}
\fancyhf{}
\rhead{PeshMind}
\lhead{Relazione Tecnica}
\rfoot{Pagina \thepage}

\title{\textbf{PeshMind} \\
\large Prolog-based Engine for System \& Host Management and INference Daemon \\
\vspace{0.5cm}
Relazione Tecnica del Progetto}
\author{Mirko Mariotti}
\date{Marzo 2026}

\begin{document}

\maketitle
\thispagestyle{empty}

\newpage
\tableofcontents
\newpage

\section{Introduzione}

PeshMind è un sistema per la scoperta e l'analisi della topologia di una rete LAN che si basa esclusivamente
sugli indirizzi MAC osservati dagli switch di rete.
Il sistema combina un'applicazione CLI scritta in Go con un motore Prolog per identificare 
le connessioni di rete, gli switch interni/di accesso e i potenziali "ghost switch" (switch non censiti)
presenti nell'infrastruttura di rete.

\subsection{Motivazione e Obiettivi}

In reti complesse di grandi dimensioni, la documentazione della topologia fisica spesso diventa obsoleta 
o incompleta. PeshMind affronta questo problema fornendo un sistema automatizzato che:

\begin{itemize}
    \item Ricostruisce la topologia di rete analizzando solo le tabelle MAC degli switch
    \item Identifica automaticamente switch non censiti ("ghost switch"), abilitando la possibilità di scoprire anomalie e dispositivi non autorizzati
    \item Distingue tra switch di core/distribuzione e switch di accesso
    \item Fornisce visualizzazioni grafiche della topologia inferita
\end{itemize}

\section{Architettura del Sistema}

\subsection{Componenti Principali}

Il sistema PeshMind è composto da tre componenti principali che interagiscono tra loro:

\subsubsection{Applicazione CLI Go}

L'applicazione CLI \textbf{peshmind}, sviluppata in Go, coordina tutte le operazioni del sistema, comprese le
due altre componenti. I suoi sottocomandi, realizzati con il framework Cobra, gestiscono le seguenti funzionalità:

\begin{itemize}
    \item Recupero automatico delle tabelle MAC dagli switch fisici
    \item Conversione dei dati in fatti Prolog
    \item Avvio e gestione del motore Prolog
    \item Interrogazione del sistema di inferenza
    \item Generazione di visualizzazioni DOT
    \item Simulazione di topologie di rete
\end{itemize}

\subsubsection{Motore di Inferenza Prolog}

Il motore Prolog (SWI-Prolog) analizza i fatti relativi agli indirizzi MAC e inferisce la topologia di rete attraverso regole logiche. Utilizza la libreria Pengines per esporre un endpoint HTTP RESTful.

\subsubsection{Knowledge Base}

La base di conoscenza memorizza i dati degli switch e degli indirizzi MAC sotto forma di fatti Prolog:

\begin{itemize}
    \item \texttt{switch(MacAddr)} - Dichiara l'esistenza di uno switch
    \item \texttt{switchname(MacAddr, Name)} - Associa un MAC a un nome
    \item \texttt{seen(SwitchMac, DeviceMac, Port)} - Registra MAC visti su porte specifiche
\end{itemize}

\subsection{Due possibili flussi di lavoro}

Il sistema supporta due flussi di lavoro principali, quello di analisi di una rete esistente e 
quello di simulazione di topologie di rete. Il secondo è particolarmente utile per testare e validare le regole
di inferenza Prolog prima di applicarle a dati reali.

\begin{multicols}{2}
        \textbf{Analisi Rete Esistente}
        \begin{enumerate}
            \item Configurazione switch in \texttt{peshmind.json}
            \item Acquisizione tabelle MAC via SSH o telnet
            \item Conversione dati in fatti Prolog
            \item Caricamento KB e inferenza
            \item Interrogazione e visualizzazione risultati
        \end{enumerate}
\columnbreak
        \textbf{Simulazione Topologia di Rete}
        \begin{enumerate}
            \item Creazione file DOT con topologia desiderata
            \item Esecuzione simulazione per generare fatti Prolog
            \item Caricamento KB e inferenza
            \item Interrogazione e visualizzazione risultati
            \item Iterazione per testare diverse configurazioni
        \end{enumerate}
\end{multicols}

\subsection{Struttura del Codice}

I sorgenti del progetto sono organizzati nella seguente struttura:

\begin{verbatim}
peshmind/
|-- main.go              # Entry point dell'applicazione
|-- cmd/                 # Comandi Cobra
|   |-- root.go         # Comando radice
|   |-- list.go         # Elenca switch configurati
|   |-- update.go       # Aggiorna dati switch
|   |-- serve.go        # Avvia motore Prolog
|   |-- query.go        # Interroga il motore
|   |-- dot.go          # Genera file DOT
|   +-- simulate.go     # Simula topologie
|-- pkg/peshmind/        # Logica di business
|   |-- peshmind.go     # Configurazione
|   |-- switch.go       # Gestione switch
|   |-- spawn.go        # Avvio motore Prolog
|   +-- simulate.go     # Logica di simulazione
+-- kbpool/              # Knowledge base Prolog
    |-- rules.pl        # Regole di inferenza
    |-- server.pl       # Server HTTP Pengines
    +-- createkb.sh     # Script creazione KB
\end{verbatim}
\subsection{Tecnologie Utilizzate}

\subsubsection{Go}

Il linguaggio Go è stato scelto per l'applicazione CLI in quanto offre performance eccellenti,
 consente la compilazione statica con generazione di un singolo binario facilmente deployable, 
 fornisce un supporto eccellente per la programmazione concorrente e mette a disposizione un ricco
  ecosistema di librerie pronte all'uso. \\
Le librerie Go che sono state utilizzate includono:
\begin{itemize}
    \item \textbf{Cobra}: Framework CLI per comandi e flag
    \item \textbf{Pengine Client}: Client Go per interrogare Pengines
    \item \textbf{go-graphviz}: Parsing e generazione di grafi DOT
    \item \textbf{golang.org/x/term}: Input sicuro password
\end{itemize}
\subsubsection{SWI-Prolog}

SWI-Prolog fornisce il motore di inferenza grazie a un potente sistema di inferenza logica, alla 
libreria Pengines per l'esposizione di API HTTP, a una sintassi dichiarativa particolarmente adatta alla 
definizione di regole complesse e a ottime prestazioni anche su knowledge base di grandi dimensioni.

\section{Utilizzo}
\subsection{Comandi CLI generali}

Gran parte dei comandi CLI di PeshMind sono utilizzabili sia per l'analisi di una rete reale che per la 
simulazione di topologie, in quanto operano sui fatti Prolog generati da entrambe le modalità.

\subsubsection{List}

Il comando \texttt{list} visualizza tutti gli switch configurati con i loro dettagli:

\begin{lstlisting}
peshmind list
\end{lstlisting}

\subsubsection{Serve}

Il comando \texttt{serve} avvia il motore Prolog con endpoint HTTP:

\begin{lstlisting}
peshmind serve --kbpool kbpool
\end{lstlisting}

Esegue le seguenti operazioni:
\begin{itemize}
    \item Verifica la presenza di SWI-Prolog
    \item Esegue \texttt{createkb.sh} per consolidare i fatti
    \item Avvia il server HTTP sulla porta 3030
    \item Rimane in esecuzione fino a Ctrl+C
\end{itemize}

\subsubsection{Query}

Il comando \texttt{query} permette di interrogare il motore Prolog:

\begin{lstlisting}
# Trova tutti gli switch
peshmind query "switch(X)"

# Trova switch direttamente connessi
peshmind query "directn(X,Y)"

# Identifica ghost switch
peshmind query "ghostn(X,Y)"
\end{lstlisting}

\subsubsection{Dot}

Il comando \texttt{dot} genera una rappresentazione grafica della topologia:

\begin{lstlisting}
peshmind dot > network.dot
dot -Tpng network.dot -o network.png
\end{lstlisting}

Il file DOT generato include:
\begin{itemize}
    \item Cluster per ogni switch con le relative porte
    \item Collegamenti fisici tra le porte
    \item Ghost switch evidenziati con stile differente
    \item Codifica colori per migliore leggibilità
\end{itemize}

\subsection{Comandi per l'analisi di una rete reale}

\subsubsection{Update}

Il comando \texttt{update} recupera le tabelle MAC dagli switch:

\begin{lstlisting}
# Aggiorna uno switch specifico
peshmind update --switch switch-01

# Aggiorna tutti gli switch
peshmind update --switch all
\end{lstlisting}

Questo comando:
\begin{enumerate}
    \item Si connette allo switch via SSH usando expect
    \item Esegue comandi per ottenere la tabella MAC
    \item Parsifica l'output con espressioni regolari
    \item Genera fatti Prolog nel formato \texttt{seen/3}
    \item Salva il file nella directory \texttt{kbpool}
\end{enumerate}

\subsection{Comandi per la simulazione di topologie}

\subsubsection{Simulate}

Il comando \texttt{simulate} permette di testare topologie:

\begin{lstlisting}
peshmind simulate --simname topology.dot \
    --emit-dot output.dot \
    --output facts.pl
\end{lstlisting}

Anch'esso. alla pari del comando \texttt{update}, genera fatti Prolog, ma invece di connettersi a switch reali,
genera i fatti in maniera programmatica a partire da un file DOT che descrive la topologia desiderata. Supporta
attributi speciali per identificare ghost switch, specificare MAC address e porte di connessione.

\section{Inferenza Logica con Prolog}

Il cuore del sistema è il motore di inferenza Prolog, che utilizza i fatti generati per dedurre la topologia di rete.
In questa sezione vengono descritte le principali regole di inferenza implementate in Prolog.

\subsection{Predicato far/2}

Due switch sono "lontani" se sono connessi a un terzo switch sulla stessa porta:

\begin{lstlisting}
far(X,Y) :- 
    seen(X,K,MIDPORTX), 
    seen(Y,K,MIDPORTY),
    seen(X,Y,PORTX),
    seen(Y,X,PORTY),
    switch(K), switch(X), switch(Y), 
    X \= Y, Y \= K, X \= K, 
    PORTX == MIDPORTX,
    PORTY == MIDPORTY.
\end{lstlisting}

Questo predicato identifica la presenza di uno switch intermedio tra due switch.

\subsection{Predicato direct/2}

Due switch sono direttamente connessi se non sono lontani:

\begin{lstlisting}
direct(X,Y) :- \+ far(X,Y).
\end{lstlisting}

\subsection{Predicato directp/4}

Connessione diretta con informazioni sulle porte:

\begin{lstlisting}
directp(X,Y,PORTX,PORTY) :- 
    seen(X,Y,PORTX),
    seen(Y,X,PORTY),
    \+ far(X,Y).
\end{lstlisting}

\subsection{Predicato ghost/2}

Identifica ghost switch tra due switch direttamente connessi:

\begin{lstlisting}
ghost(X,Y) :- 
    directp(X,Y,PORTX,PORTY),
    seen(X,Z,PORTX),
    seen(Y,Z,PORTY),
    \+ switch(Z).
\end{lstlisting}

Se due switch vedono lo stesso MAC (non switch) sulle porte di interconnessione, esiste un ghost switch.

\subsection{Predicato internalswitch/1}

Identifica switch interni (core/distribuzione):

\begin{lstlisting}
internalswitch(X) :- 
    seen(X,Y,PORT1),
    seen(X,Z,PORT2),
    switch(X), switch(Y), switch(Z), 
    X \= Y, X \= Z, Y \= Z,
    PORT1 \= PORT2.
\end{lstlisting}

Uno switch è interno se è connesso ad almeno due switch diversi.

\subsection{Predicato edgeswitch/1}

Identifica switch di accesso:

\begin{lstlisting}
edgeswitch(X) :- switch(X), \+ internalswitch(X).
\end{lstlisting}

\section{Dettagli di Configurazione}

\subsection{File di Configurazione}

Il file \texttt{peshmind.json} definisce gli switch e i parametri di sistema in entrambe le modalità 
(di analisi e simulazione). Ecco un esempio di configurazione:

\begin{lstlisting}
{
  "debug": true,
  "endpoint": "http://localhost:3030/pengine",
  "switches": {
    "switch-01": {
      "name": "switch-01",
      "mac": "aabbccddeeff",
      "description": "Core Switch Piano 1",
      "ip": "192.168.1.10",
      "username": "admin",
      "password": "ask",
      "model": "HP-Aruba",
      "port": 48,
      "extraports": [],
      "data": ""
    }
  }
}
\end{lstlisting}

\begin{itemize}
    \item \textbf{debug}: Abilita output verboso
    \item \textbf{endpoint}: URL dell'endpoint Pengines HTTP
    \item \textbf{switches}: Dizionario di configurazioni switch
    \begin{itemize}
        \item \textbf{name}: Identificatore univoco dello switch
        \item \textbf{mac}: Indirizzo MAC dello switch (senza separatori)
        \item \textbf{description}: Descrizione leggibile
        \item \textbf{ip}: Indirizzo IP di gestione
        \item \textbf{username}: Nome utente SSH
        \item \textbf{password}: Password SSH (usa "ask" per input interattivo)
        \item \textbf{model}: Modello switch (HP-Aruba, HP-Aruba-Press)
        \item \textbf{port}: Numero di porte fisiche
        \item \textbf{extraports}: Porte aggiuntive da monitorare
    \end{itemize}
\end{itemize}

\subsection{Modalità Analisi Rete Esistente}
Per analizzare una rete esistente, è necessario solamente configurare correttamente il file \texttt{peshmind.json} 
con le informazioni di accesso per ogni switch gestito.

\subsection{Modalità Simulazione}

Nel caso della simulazione, in aggiunta alla configurazione, è necessario creare un file DOT che descriva la topologia
di rete desiderata. Il file DOT deve seguire una struttura specifica per essere analizzato correttamente dal sistema
di simulazione.
La modalità di simulazione permette di testare topologie usando file DOT. Il sistema:

\begin{enumerate}
    \item Parsifica il grafo DOT usando la libreria \texttt{go-graphviz}
    \item Identifica nodi switch (normali e ghost)
    \item Estrae le connessioni dagli archi del grafo
    \item Genera MAC address fittizi per dispositivi end-node
    \item Produce fatti Prolog equivalenti
    \item Supporta attributi speciali:
    \begin{itemize}
        \item \texttt{ghost="true"} - Marca ghost switch
        \item \texttt{mac="..."} - Specifica MAC per ghost switch
        \item \texttt{port="N"} - Specifica porta di connessione
        \item \texttt{ports="N"} - Numero di porte per ghost switch
    \end{itemize}
\end{enumerate}

\section{Esempi di Utilizzo}

\subsection{Scoperta Topologia di Rete Esistente}

\textbf{Scenario}: Un amministratore di rete deve documentare una rete esistente con documentazione obsoleta.

\textbf{Procedimento}:
\begin{enumerate}
    \item Configurare tutti gli switch gestiti in \texttt{peshmind.json}
    \item Eseguire \texttt{peshmind update --switch all}
    \item Avviare il motore: \texttt{peshmind serve}
    \item Interrogare: \texttt{peshmind query "directpn(X,Y,PX,PY)"}
    \item Generare visualizzazione: \texttt{peshmind dot > topology.dot}
\end{enumerate}

\subsection{Identificazione Ghost Switch}

\textbf{Scenario}: Identificare switch non gestiti o non documentati nella rete.

\textbf{Procedimento}:
\begin{enumerate}
    \item Acquisire dati da switch gestiti
    \item Interrogare: \texttt{peshmind query "ghostn(X,Y)"}
    \item Analizzare i risultati per identificare switch mancanti
    \item Aggiungere gli switch alla configurazione se necessario
\end{enumerate}

\subsection{Classificazione Switch}

\textbf{Scenario}: Determinare quali switch sono core/distribuzione e quali sono di accesso.

\textbf{Procedimento}:
\begin{enumerate}
    \item Switch interni: \texttt{peshmind query "internalswitchn(X)"}
    \item Switch di accesso: \texttt{peshmind query "edgeswitchn(X)"}
    \item Analizzare la distribuzione per ottimizzare la rete
\end{enumerate}

\subsection{Testing Modifiche alla Rete}

\textbf{Scenario}: Validare una nuova topologia prima dell'implementazione fisica.

\textbf{Procedimento}:
\begin{enumerate}
    \item Creare un file DOT della topologia proposta
    \item Eseguire simulazione: \texttt{peshmind simulate --simname new-topology.dot}
    \item Verificare connettività e relazioni
    \item Identificare potenziali problemi prima dell'implementazione
\end{enumerate}



\section{Limitazioni e Sviluppi Futuri}

\subsection{Limitazioni Attuali}

\begin{itemize}
    \item Supporto limitato a modelli HP-Aruba
    \item Richiede accesso SSH agli switch
    \item Non gestisce VLAN multiple in modo esplicito
    \item Visualizzazioni DOT richiedono post-processing manuale
\end{itemize}

\subsection{Sviluppi Futuri}

\begin{enumerate}
    \item \textbf{Supporto Multi-Vendor}
    \begin{itemize}
        \item Aggiungere template per Cisco, Juniper, Arista
        \item Supporto SNMP come alternativa a SSH
    \end{itemize}
    
    \item \textbf{Analisi VLAN}
    \begin{itemize}
        \item Inferire configurazioni VLAN
        \item Identificare errori di configurazione VLAN
    \end{itemize}
    
    \item \textbf{Interface Web}
    \begin{itemize}
        \item Dashboard per visualizzazione interattiva
        \item Editor grafico topologie
        \item Monitoring real-time
    \end{itemize}
    
    \item \textbf{Analisi Avanzate}
    \begin{itemize}
        \item Identificazione loop
        \item Analisi percorsi ridondanti
        \item Suggerimenti ottimizzazione
    \end{itemize}
    
    \item \textbf{Integrazione}
    \begin{itemize}
        \item Export verso NetBox, LibreNMS
        \item API REST per integrazione
        \item Webhook per notifiche cambiamenti
    \end{itemize}
\end{enumerate}

\section{Conclusioni}

PeshMind rappresenta un approccio innovativo alla scoperta automatica di topologie di rete, combinando l'efficienza di Go con la potenza dell'inferenza logica Prolog. Il sistema si è dimostrato efficace nell'identificare automaticamente:

\begin{itemize}
    \item Connessioni fisiche tra switch
    \item Switch non documentati (ghost switch)
    \item Classificazione switch (core vs. accesso)
    \item Anomalie nella topologia di rete
\end{itemize}

La combinazione di un'interfaccia CLI intuitiva con un motore di ragionamento logico fornisce agli amministratori di rete uno strumento potente per comprendere, documentare e validare le loro infrastrutture di rete.

Il progetto è rilasciato con licenza Apache 2.0, favorendo l'adozione e la contribuzione della community open source.

\section{Bibliografia e Riferimenti}

\begin{itemize}
    \item SWI-Prolog Documentation: \url{https://www.swi-prolog.org/}
    \item Pengines Library: \url{https://www.swi-prolog.org/pldoc/doc_for?object=section(%27packages/pengines.html%27)}
    \item Cobra CLI Framework: \url{https://github.com/spf13/cobra}
    \item Graphviz Documentation: \url{https://graphviz.org/}
    \item Go Programming Language: \url{https://golang.org/}
\end{itemize}

\appendix

\end{document}
